\section{The original stationary microscopic local limit}
\label{sec:stationary-microscopic-LLT}
We now evaluate the complete singleton physical probability. The
initial and terminal cell offsets in
\eqref{eq:exact-physical-cell-correction} remain present throughout the
argument. The overlap kernel is tested against the local endpoint
measure, rather than estimated through a growing Fourier norm.
This proves the disappearance of the signed complement in
\eqref{eq:stationary-middle-complement}.

Fix once and for all the convex polygonal fundamental cell
$\mathcal P$ of Section~\ref{sec:exact-stationary-endpoint}. The finite
set $\mathcal D$ and $J_*=|\mathcal D|$ depend only on this cell, the
lattice and the uniform horizon. They are not enlarged with $m,t$ or
with a frequency cutoff.

\begin{lemma}[Admissibility of the actual overlap tests]
\label{lem:overlap-local-admissibility}
For $d,e\in\mathcal D$ the functions
\[
 F_R^{d,e}(x,s,y)=H_{I_{R,d}(x),I_{R,e}(y)}(-s)
\]
have a common bound and support $|s|\le\tau_+$. They satisfy the
fixed and parameter-varying test conditions of
Corollary~\ref{cor:local-endpoint-measures}. Moreover
\begin{equation}\label{eq:overlap-total-amplitude}
 \sum_{d,e}\int F_R^{d,e}(x,s,y)\,d\nu(x)\,ds\,d\nu(y)
                                      =\bar\tau_R^2.
\end{equation}
\end{lemma}
\begin{proof}
An overlap of two flight intervals is bounded by $\tau_+$ and vanishes
outside $[-\tau_+,\tau_+]$. At a regular collision state, excluding
flights parallel to or passing through a vertex of a cell edge, the
finite interval endpoints are given by intersections of a straight
line with finitely many fixed lines, truncated by the regular entry
root. They depend continuously on the parameter and collision state.
Taking maxima and minima and allowing empty intervals preserves
continuity of the intervals in symmetric-difference length. Since
\[
 |H_{I,J}(s)-H_{I',J'}(s)|
                   \le |I\mathbin\triangle I'|+|J\mathbin\triangle J'|,
\]
this implies uniform-in-$s$ continuity of the overlap.

The exceptional sets just described, together with grazing and the
one-flight singularity set, have zero collision measure: they are
finitely many piecewise analytic curves and boundary pieces, with
any coincident edge segment assigned the fixed cell convention. A
flight contained in an edge has a prescribed position-direction
constraint and also has zero collision measure. On every compact set
separated from these sets at $R_0$, the selected roots and all
transverse crossings persist uniformly for $R$ near $R_0$. Thus the
interval endpoints, and hence the overlap tests, converge uniformly
there. Open neighborhoods of the exceptional sets can be chosen with
arbitrarily small $\nu$ measure and null boundary. The union of their
initial and terminal lifts to a compact $s$ slab has arbitrarily small
product measure. This proves the varying-test condition; it does not
rely on an unquantified exchange of two pointwise limits.

Finally $\int_\R H_{I,J}(s)\,ds=|I||J|$. Sum over the two finite
flight partitions to obtain $\tau_R(x)\tau_R(y)$ and integrate the
independent endpoint measures. This gives
\eqref{eq:overlap-total-amplitude}.
\end{proof}

\begin{theorem}[Uniform local limit for the original stationary event]
\label{thm:stationary-microscopic-LLT}
For every $M<\infty$, uniformly for $R\in I$, $k\in\Z^2$, $m\ge1$
and $t>0$ in the indicated central sets,
\begin{align}
 m^{3/2}\mathbb P_R\{W_R(t)=k,C_R(t)=m\}
 &=\bar\tau_Rg_{\Sigma_R}(\zeta_{m,R}(k,t))+o(1),
           &&|\zeta_{m,R}(k,t)|\le M,
                                      \label{eq:stationary-collision-scale-LLT}\\
 t^{3/2}\mathbb P_R\{W_R(t)=k,C_R(t)=m\}
 &=g_{\mathcal V_R}(\xi_{t,R}(k,m))+o(1),
           &&|\xi_{t,R}(k,m)|\le M,
                                      \label{eq:stationary-physical-scale-LLT}
\end{align}
where
\[
 \xi_{t,R}(k,m)=\left(\frac{k}{\sqrt t},
                       \frac{m-t/\bar\tau_R}{\sqrt t}\right),
 \qquad
 \mathcal V_R=\bar\tau_R^{-1}
 D_R^{\rm clk}\Sigma_R(D_R^{\rm clk})^{\mathsf T},\quad
 D_R^{\rm clk}=\operatorname{diag}(1,1,-1/\bar\tau_R).
\]
The first error tends to zero as $m\to\infty$ and the second as
$t\to\infty$. In particular there is $c_M>0$ such that, for all
sufficiently large $t$ in the second central set,
\begin{equation}\label{eq:stationary-microscopic-denominator}
 \mathbb P_R\{W_R(t)=k,C_R(t)=m\}\ge c_M t^{-3/2}.
\end{equation}
\end{theorem}
\begin{proof}
Use the exact overlap formula
\eqref{eq:exact-overlap-probability}. Its $(d,e)$ term has the exact
collision label $k+d-e$, not $k$. Multiplying that term by $m^{3/2}$
expresses it as $\bar\tau_R^{-1}\rho_{m,R}^{k+d-e,t}(F_R^{d,e})$.
The offsets belong to a fixed finite set, so
$\zeta_{m,R}(k+d-e,t)-\zeta_{m,R}(k,t)=O(m^{-1/2})$.

For any sequence in the first central set choose a subsequence with
$R\to R_0$ and $\zeta\to\zeta_0$. Apply
Corollary~\ref{cor:local-endpoint-measures} and
Lemma~\ref{lem:overlap-local-admissibility} to every one of the
finitely many pairs. The sum of their limiting integrals is
$\bar\tau_{R_0}^{-1}g_{\Sigma_{R_0}}(\zeta_0)\bar\tau_{R_0}^2$.
This proves the first limit by the sequential criterion for uniformity.
Both cell corrections and the initial stationary age were retained
in every term.

In the physical central set $t/m=\bar\tau_R+O(t^{-1/2})$, uniformly,
and the two central sets are comparable. The matrix identity
\eqref{eq:physical-covariance-jacobian} gives
$\det\mathcal V_R=\det\Sigma_R/\bar\tau_R^5$.
Equivalently, for
$A_R=\sqrt{\bar\tau_R}\operatorname{diag}(1,1,-\bar\tau_R)$,
\[
 A_R\mathcal V_RA_R^{\mathsf T}=\Sigma_R,\qquad
 |\det A_R|=\bar\tau_R^{5/2},\qquad
 \zeta_{m,R}=A_R\xi_{t,R}+O(t^{-1/2}).
\]
Gaussian change of variables converts the first formula into the
second. Uniform ellipticity, continuity and compactness imply
$\inf_{R,|\xi|\le M}g_{\mathcal V_R}(\xi)>0$. The uniform error
then proves \eqref{eq:stationary-microscopic-denominator}.
\end{proof}

\begin{corollary}[The full signed middle complement vanishes]
\label{cor:stationary-middle-vanishes}
For every fixed $P>0$ take the explicit bandwidth $B_m=m^{P+3/2}$
and the physical central radius $r_m=2m^{-99/200}$. For the \emph{whole}
complement $\mathfrak M_{m,R,B_m}$ defined in
\eqref{eq:stationary-middle-complement},
\begin{equation}\label{eq:stationary-middle-vanishes}
 \sup_{\substack{R\in I,\ k\in\Z^2,\ t>0\\
                         |\xi_{t,R}(k,m)|\le M}}
 t^{3/2}|\mathfrak M_{m,R,B_m}(k,t)|\longrightarrow0.
\end{equation}
In particular $B_m=m^{5/2}$ is one admissible polynomial bandwidth.
\end{corollary}
\begin{proof}
Theorem~\ref{thm:physical-singleton-reduction} gives the exact
Gaussian-plus-complement identity with central error
$O(m^{-11/700}\sqrt{\log(2+m)})$ and normalized far error $O(m^{-P})$.
Theorem~\ref{thm:stationary-microscopic-LLT} independently evaluates
its full probability, with the same Gaussian and normalization.
Subtract the two formulas. This proves
\eqref{eq:stationary-middle-vanishes} on precisely the stated domain.

The mechanism is fixed-band spectral decay followed by positive
interval envelopes and the exact endpoint overlap. The complement
contains the entire noncentral spatial torus and every roof frequency
between the central ball and $B_m$; none has been omitted. This proves
a signed inverse estimate, not an absolute integral estimate for the
modulus of its transform. No uniform-in-$B$ spectral constant, or
unproved cancellation assumption, is used in the subtraction.
\end{proof}

\subsection{Endpoint selectors and normalization on the same trajectory}
Let $\chi_R,\psi_R$ be bounded functions of a collision state and its
age in the current flight. For definiteness assume they are restrictions
of a jointly continuous bounded family on
$I\times\overline M\times[0,\tau_+]$. The same statements hold for
indicators of fixed endpoint sets whose boundaries have zero stationary
measure, with the uniform exceptional-neighborhood condition used in
Corollary~\ref{cor:local-endpoint-measures}. Write
\[
 \mu_R(\chi_R)=\bar\tau_R^{-1}\int_M\int_0^{\tau_R(x)}
                       \chi_R(x,a)\,da\,d\nu(x).
\]
On a trajectory starting at $(x,a)$ let $(y_t,v_t)$ be its current
collision state and flight age at time $t$.

\begin{theorem}[A microscopic local law with endpoint selectors]
\label{thm:stationary-endpoint-selected-LLT}
For the endpoint class just specified, uniformly on fixed physical
central compact sets,
\begin{align}
 t^{3/2}\mathbb E_R[\chi_R(x,a)\psi_R(y_t,v_t)
             \one_{\{W_R(t)=k,C_R(t)=m\}}]
 =\mu_R(\chi_R)\mu_R(\psi_R)
            g_{\mathcal V_R}(\xi_{t,R}(k,m))+o(1).
 \label{eq:stationary-endpoint-selected-LLT}
\end{align}
Thus the conditional distribution of the initial and terminal flight
states, tested against this endpoint class, tends to the product of
the two stationary flight-state laws. For nonnegative selectors with
$\inf_R\mu_R(\chi_R)\mu_R(\psi_R)>0$, their selected microscopic
mass has a uniform positive lower bound of order $t^{-3/2}$.
\end{theorem}
\begin{proof}
In the $(d,e)$ term of the exact source formula replace the overlap by
\[
 H_R^{d,e;\chi,\psi}(x,y,s)=
 \int_{I_{R,d}(x)}\chi_R(x,a)
        \one_{I_{R,e}(y)}(s+a)\psi_R(y,s+a)\,da.
\]
The same common support and bounds hold. Continuity of the selectors,
the interval symmetric-difference estimate and dominated integration
give the admissibility conditions of
Lemma~\ref{lem:overlap-local-admissibility}. For the stated indicator
extension, apply the null-boundary approximation on the product of
endpoint state--age spaces before integrating the ages. Fubini makes
the exceptional overlap tests null for $\nu(dx)ds\nu(dy)$, and the
specified neighborhood condition gives parameter uniformity.

Integration in $s$ changes variables $v=s+a$ and factors the amplitude:
\begin{align*}
 &\sum_{d,e}\int H_R^{d,e;\chi,\psi}(x,y,s)
                           \,d\nu(x)\,ds\,d\nu(y)\\
 &\quad=\left(\int_M\int_0^{\tau_R(x)}\chi_R(x,a)\,da\,d\nu(x)\right)
   \left(\int_M\int_0^{\tau_R(y)}\psi_R(y,v)\,dv\,d\nu(y)\right).
\end{align*}
Repeat the preceding local-measure proof and the same covariance change
of variables. This gives \eqref{eq:stationary-endpoint-selected-LLT}.
Division by \eqref{eq:stationary-microscopic-denominator} proves the
conditional assertion. Positivity of the two stationary means and the
Gaussian gives the selected denominator bound.
\end{proof}

\begin{corollary}[Normalized polynomial-band reconstruction]
\label{cor:normalized-physical-band}
Let $A=\{W_R(t)=k,C_R(t)=m\}$ be in a fixed physical central compact
set, and let $\lambda_{A,B}$ be the positive likelihood in
\eqref{eq:physical-time-positive-likelihood}. For $B_m=m^{P+3/2}$,
\begin{equation}\label{eq:normalized-physical-band}
 \left\|\frac{\one_A\,d\mathbb P_R}{\mathbb P_R(A)}
       -\frac{\lambda_{A,B_m}\,d\mathbb P_R}
                           {\mathbb E_R\lambda_{A,B_m}}\right\|_{\TV}
                            =O(m^{-P}).
\end{equation}
The comparison is on the original trajectory space and is uniform
against every bounded measurable trajectory selector. It uses no event
replacement. The same comparison holds after multiplication by fixed
nonnegative endpoint selectors from
Theorem~\ref{thm:stationary-endpoint-selected-LLT} with uniformly
positive selected mass on the normalized scale.
\end{corollary}
\begin{proof}
Theorem~\ref{thm:physical-time-selector-reconstruction} bounds the
unnormalized source error by $C/B$. For two finite positive measures
$\alpha,\beta$ with masses $p,q>0$,
\[
 \|\alpha/p-\beta/q\|_{\TV}
 \le\frac{\|\alpha-\beta\|_{\TV}+|q-p|}{p}
 \le\frac{2\|\alpha-\beta\|_{\TV}}p.
\]
Use the newly proved denominator $p\ge c_M t^{-3/2}$, and note that
$t/m$ stays in a fixed positive compact interval. For large $m$, the
source error is smaller than $p/2$, so $q>0$. This proves
\eqref{eq:normalized-physical-band}. Multiplication by a bounded
nonnegative endpoint weight contracts the unnormalized error up to
its supremum norm; the selected denominator is provided by the preceding
theorem. Total variation controls every bounded test on the unchanged
trajectory space, without any regularity of that additional test.
\end{proof}

\paragraph{The two local problems.}
Theorems~\ref{thm:stationary-microscopic-LLT} and
\ref{thm:stationary-endpoint-selected-LLT} prove local statements for
the original stationary event at prescribed physical collision count.
They do not identify a Gaussian bridge for arbitrary long path selectors.
The arbitrary-selector assertion in
Corollary~\ref{cor:normalized-physical-band} is a comparison of two
posterior measures, not a formula for the selected Gaussian amplitude.
Nor does the positive age-overlap mechanism create a pointwise density
estimate for the four-coordinate return record $J_{n,R}$. The latter's
common raw correction and return-frequency complement retain their
separate roles in Theorem~\ref{thm:LLT}. The original raw return problem
and all its geometric and probabilistic inputs remain in the article.
